% T22 -- CAN/CNA, num\'erisation et \'echantillonnage -- banque d'exercices
%
% COMPL\'EMENTARIT\'E AVEC LA FICHE fiches/T22_CAN_CNA_Numerisation.tex
% Les trois Questions Types de la fiche sont :
%   QT1 = 2013 E1 (tour de laboratoire), questions 42 et 43 : plage -50/+50 degC,
%         pr\'ecision 0,1 degC, 0-10 V, nombre de bits + mot hexa 732_16 pour theta = -5 degC.
%         (\'Enonc\'e authentique : 2013/E1_Sujet_2013_commun_SII.md l.507-509.)
%   QT2 = 2026 E2 (BAC+3), question B11 : sensibilit\'e S = 2,684 LSB/degC du CAN 8 bits
%         ADS7029-Q1 et mot num\'erique N_C pour T_C = 55 degC.
%         C'est la QT explicitement signal\'ee par le CLAUDE.md du projet comme ayant \'et\'e
%         corrig\'ee lors de l'audit 2026-08-24 : elle N'EST PAS reprise ici.
%   QT3 = 2026 E2 (BAC+3), question B17 : th\'eor\`eme de Shannon + r\^ole et fr\'equence de
%         coupure du filtre anti-repliement.
% AUCUN des exercices ci-dessous ne refait ce travail intellectuel :
%   - le sujet 2013 n'est pas utilis\'e du tout ;
%   - du sujet 2026 on n'exploite QUE les questions B10, B12 (lecture datasheet + erreur de
%     quantification) et B13/B15/B16 (p\'eriode d'\'echantillonnage, E/B, temps de conversion) ;
%     la sensibilit\'e S et l'\'enonc\'e de Shannon ne sont jamais redemand\'es ;
%   - le raisonnement "Shannon -> f_c = f_e/2" de la QT3 n'appara\^it nulle part : l'exercice 4
%     porte sur le crit\`ere de dynamique de boucle (10 \'echantillons dans le r\'egime
%     transitoire), qui est un crit\`ere DIFF\'ERENT et ind\'ependant de Shannon.
%
% SOURCES (toutes relues dans les .md de conversion Marker, chemin racine
%  docs/CAPET_Markdown-20260629T134311Z-3-001/CAPET_Markdown/docs/CAPET_SII_IE/06_PDFs_Originaux) :
%
%   - Ex. 1 : 2022, "\'EPREUVE \'ECRITE DISCIPLINAIRE" (E2_Disciplinaire_IE_2022.md l.20 pour
%     le libell\'e exact de l'\'epreuve), questions 23 et 24.
%     \'Enonc\'e : l.453-457 (PIC, ADC0/ADC1/ADC2, seuil diodes 0,6 V), l.455 ("La tension de
%     r\'ef\'erence du Convertisseur Analogique Num\'erique est \'egale \`a 3,5 V"),
%     l.458-468 (questions 23 et 24 mot pour mot), l.522 ("CAN 12 bits et sa tension de
%     r\'ef\'erence de 3,5V" -- confirme la r\'esolution, le document DT D3 \'etant une image).
%     CORRIG\'E OFFICIEL : Rapport_Jury_2022.md l.303-319 :
%       * "C'est un convertisseur 12 bits", 16 entr\'ees AN0..AN15, multiplexage par CHS3:CHS0 ;
%       * "Quantum = (vREF+ - vREF-)/(2^12 - 1) = (3,5-0)/4095 = 0,85 mV / LSB" ;
%       * U_Peltier = 2 V -> ADC2 = 1 V (pont diviseur R4/R5) -> N = 1/q = 1170 ;
%       * N = 175 -> ADC2 = 175*3,5/4095 = 0,15 V -> U_Peltier = 0,3 V, "donc inf\'erieure au
%         0,4V minimum impos\'e sur le diagramme des exigences".
%     RECALCUL\'E : 3,5/4095 = 0,854 7 mV ; 1/0,000 854 7 = 1 170,0 ; 175*3,5/4095 = 0,149 6 V ;
%     x2 = 0,299 V. Tout concorde avec le rapport.
%     CONVENTION : le corrig\'e officiel 2022 utilise explicitement 2^n - 1.
%
%   - Ex. 2 : 2023, "\'EPREUVE \'ECRITE DSICIPLINAIRE" [sic, coquille du sujet lui-m\^eme,
%     E1_Sujet_2023.md l.11], questions D2 et D3. \'Enonc\'e : l.361-369 (codeur incr\'emental
%     0,15 mm/impulsion, excursion 680 mm, compteur \`a 0 v\'erin rentr\'e, pr\'ecision requise
%     5 mm, "ne retenant que les 8 bits de poids fort").
%     PAS de corrig\'e officiel disponible (aucun Rapport_Jury_2023 dans le corpus) :
%     TOUT EST RECALCUL\'E ICI, deux fois.
%       680/0,15 = 4 533,3 -> code max entier 4 533 ; 2^12 = 4 096 < 4 534 donc n = 13 bits ;
%       4 533 = 1 0001 1011 0101_2 ; garder les 8 MSB d'un mot de 13 bits = supprimer 5 LSB,
%       soit un facteur 2^5 = 32 -> pr\'ecision 0,15 x 32 = 4,8 mm < 5 mm (exigence satisfaite) ;
%       octet = 4533 >> 5 = 141 = 1000 1101_2 = 8D_16 ; 141 x 32 x 0,15 = 676,8 mm.
%
%   - Ex. 3 (docTech) : 2026, "ADMISSIBILIT\'E : \'EPREUVE 2" du CAPET (BAC + 3), option
%     Ing\'enierie \'electrique (E2_Sujet_2026_BAC3.md l.16-20 pour le libell\'e exact),
%     questions B10 et B12. \'Enonc\'e : l.417 (CAN ADS7029-Q1, documentation DTB4),
%     l.419-425 (question B10 : grandeur pleine \'echelle, r\'esolution, quantum),
%     l.429 (question B12 : erreur de quantification maximale -> erreur de mesure de T),
%     l.395 (plage de tension de 3 V, v_A = c.T), l.339-345 (cahier des charges :
%     \'etendue -5 degC \`a 90 degC, erreur de mesure < 1 degC),
%     l.636-644 (DTB5 : courbe en escalier, quantification centr\'ee, erreur en dents de scie
%     "comprise entre - q/2 et q/2").
%     ATTENTION : le document technique DTB4 est une IMAGE dans le PDF source, son texte n'est
%     pas dans la conversion Marker. L'encadr\'e docTech de l'exercice 3 est donc reconstitu\'e
%     \`a partir de la DATASHEET CONSTRUCTEUR R\'EELLE, v\'erifi\'ee par t\'el\'echargement :
%     Texas Instruments SBAS811 (janvier 2017), "ADS7029-Q1 Small-Size, Low-Power, 8-Bit,
%     2-MSPS, SAR ADC", https://www.ti.com/lit/ds/symlink/ads7029-q1.pdf
%       * titre p.1 : "8-Bit, 2-MSPS, SAR ADC" ; features : "+/-0.2 LSB DNL; +/-0.25 LSB INL" ;
%       * \S6.5 Electrical Characteristics : Resolution 8 Bits ; Full-scale input voltage span
%         0 \`a AVDD ; AVDD 2,35 V \`a 3,6 V ; INL typ +/-0,25 LSB max +/-0,5 LSB ;
%         DNL typ +/-0,2 LSB max +/-0,4 LSB ; offset error +/-0,5 LSB ; gain error +/-0,2 %FS ;
%         t_ACQ 120 ns ; throughput 2 MHz ;
%       * \S6.7 Switching Characteristics : throughput 2 MSPS, cycle time 0,5 us ;
%       * \S8.3.3 ADC Transfer Function : "1 LSB = VREF / 2^N", VREF = AVDD, N = 8.
%     -> LA DATASHEET ELLE-M\^EME RETIENT LA CONVENTION 2^n. C'est celle qu'on adopte dans
%     l'exercice 3, en signalant explicitement l'\'ecart avec 2^n - 1 (et donc avec la valeur
%     S = 2,684 LSB/degC de la question B11, qui n'est pas redemand\'ee ici).
%     Aucun Rapport_Jury_2026 n'existe dans le corpus : pas de corrig\'e officiel, tout recalcul\'e.
%
%   - Ex. 4 : 2026, m\^eme \'epreuve, questions B13, B14, B15 et B16. \'Enonc\'e : l.447-451
%     (syst\`eme du 1er ordre tau = 0,9 s, K = 4 degC/V, capteur K_c = 0,0316 V/degC),
%     l.458-472 ("Le constructeur d\'esire obtenir au moins 10 \'echantillons au cours du
%     r\'egime transitoire", B13 ; B14 r\^ole de l'E/B + sch\'ema d'un bloqueur d'ordre z\'ero ;
%     "la p\'eriode d'\'echantillonnage retenue par la suite est donn\'ee par T2 = 0,25 s" ;
%     B15 allures de u_R*(t) et u_RB(t) ; B16 temps de conversion maximal).
%     Le document r\'eponse DRB1 (l.694-700) gradue u_R de 0 \`a 0,126 4 V par pas de
%     0,031 6 V sur 0 \`a 4,5 s -- coh\'erent avec un premier ordre de gain statique 0,126 4 V
%     (= 4 degC x 0,031 6 V/degC) et de constante de temps 0,9 s.
%     RECALCUL\'E : t_R5% = 3.tau = 2,7 s ; T_e <= 2,7/10 = 0,27 s ; f_min = 1/0,27 = 3,7 Hz ;
%     avec T_2 = 0,25 s, T_conv <= 0,25 s ; le CAN convertit en 0,5 us (2 MSPS), soit un
%     rapport de 5.10^5.
%     \'Echantillons du chronogramme (u_R(k.T_2) = 0,1264.(1 - exp(-k.0,25/0,9)), en mV) :
%       k=1..11 -> 30,7 / 53,9 / 71,5 / 84,8 / 94,9 / 102,5 / 108,3 / 112,7 / 116,0 / 118,5 / 120,5
%     (une premi\`ere r\'edaction portait 32,0 / 56,3 / 74,7 ... : valeurs FAUSSES, correspondant
%      \`a tau = 0,8 s. Corrig\'ees dans le tableau ET dans les coordonn\'ees du trac\'e TikZ.)
%
%   - Ex. 5 : 2025, "\'EPREUVE \'ECRITE DISCIPLINAIRE APPLIQU\'EE" (E2_Appliquee_Commune_2025.md
%     l.19 pour le libell\'e exact), questions 23, 26 et 27. \'Enonc\'e : l.376 (d\'ebattement
%     288 deg, commandes cod\'ees sur 24 bits, "Montrer que la pr\'ecision du positionnement
%     angulaire du bras delta-theta est \'egale \`a 1.72.10^-5 deg"), l.390-394 (tension image
%     -5 V / +5 V, d\'efinition de delta-v, "Le capteur est consid\'er\'e comme suffisamment
%     pr\'ecis si son erreur de quantification q est 10 fois plus petite que la variation de
%     tension entre deux positions angulaires commandables du moteur delta-v", question 27).
%     CORRIG\'E OFFICIEL : Rapport_Jury_2025.md l.680-688 :
%       * Q23 : "delta-theta = 288/2^24 = 1.72.10^-5 deg" ;
%       * Q27 : "q = delta-v/10 ~= 6.10^-8 V. On peut d\'ecouper l'intervalle -5V,5V en N
%         segments de largeur q o\`u N v\'erifie N = 10/q ~= 1.6.10^8 ce que l'on peut coder sur
%         Nbits = [log2(1.6.10^8)] arrondi \`a l'entier sup\'erieur soit 28 bits."
%     RECALCUL\'E : 288/2^24 = 1,716 6.10^-5 deg ; delta-v = 10/2^24 = 5,960 5.10^-7 V ;
%     q = 5,960 5.10^-8 V ; N = 10/q = 1,677 7.10^8 ; log2 = 27,32 -> 28 bits. Tout concorde.
%     CONVENTION : le corrig\'e officiel 2025 utilise 2^n (N segments de largeur q sur la plage).
%
%   - Ex. 6 : EXERCICE IN\'EDIT (CNA R-2R). Aucune r\'ef\'erence de sujet ne lui est attribu\'ee.
%     Les seuls faits math\'ematiques utilis\'es (loi de sortie d'un R-2R en contre-r\'eaction,
%     erreur de mono\-tonicit\'e li\'ee \`a la tol\'erance des r\'esistances) sont recalcul\'es
%     int\'egralement dans le corrig\'e.
%
%   - Citation jury (exercice 3) : verbatim EXACT de Rapport_Jury_2020.md l.2386 :
%     "Le jury d\'eplore que la majeure partie des candidats ne sache pas d\'eterminer la
%     r\'esolution d'un capteur."
%     (La fiche T22 cite une version diff\'erente, augment\'ee d'une phrase sur l'encadrement des
%     r\'esultats qui ne figure PAS \`a cet endroit du rapport : on s'en tient ici au verbatim.)
%   - Citation jury (exercice 5) : verbatim EXACT de Rapport_Jury_2022.md l.837 (extrait) :
%     "Le jury recommande aux candidats d'apporter un soin particulier aux questions de
%     conclusion de chacune des parties. Les \'ecarts \'evalu\'es doivent \^etre clairement mis
%     en \'evidence et comment\'es."
%
% CONVENTION DE QUANTUM RETENUE DANS CE FICHIER
%   Il n'existe pas de convention unique dans les sujets CAPET eux-m\^emes :
%     - 2022 (corrig\'e officiel du jury)     -> q = V_ref / (2^n - 1)
%     - 2023 (formule impos\'ee par l'\'enonc\'e) -> q = V_ref / 2^n
%     - 2025 (corrig\'e officiel du jury)     -> q = plage / 2^n
%     - datasheet ADS7029-Q1 (\S8.3.3)        -> 1 LSB = V_REF / 2^N
%   R\`egle appliqu\'ee ici : ON SUIT LA CONVENTION DE LA SOURCE, et chaque corrig\'e l'\'enonce
%   explicitement avant de calculer. L'exercice 3 compare chiffre \`a chiffre les deux
%   conventions pour montrer que l'\'ecart (0,4 % \`a n = 8) est n\'egligeable devant les
%   erreurs de lin\'earit\'e du composant -- ce qui est l'argument attendu au CAPET.
\expandafter\ifx\csname CAPETINCLUDEMODE\endcsname\relax
\documentclass[11pt,a4paper]{article}
% preamble.tex — packages et commandes communes a toutes les fiches CAPET SII IE
\input{glyphtounicode}
\pdfgentounicode=1
\usepackage{cmap}
\usepackage[T1]{fontenc}
\usepackage[utf8]{inputenc}
\usepackage[french]{babel}
% babel[french] active a \begin{document} des raccourcis de ponctuation (:;!?)
% qui cassent la syntaxe to[...] de circuitikz (erreur "+ or - expected") ;
% on les desactive, ce qui ne change que l'espacement fin avant ces signes.
\AtBeginDocument{\shorthandoff{:;!?}}
\usepackage[a4paper, top=1.6cm, bottom=1.6cm, left=1.4cm, right=1.4cm]{geometry}
\setlength{\headheight}{14.5pt}
\usepackage{amsmath, amssymb, amsfonts}
\usepackage{siunitx}
\sisetup{locale=FR, output-decimal-marker={,}, inter-unit-product=\ensuremath{{}\cdot{}},
         group-minimum-digits=4}
% Unites non fournies par siunitx, declarees ici une fois pour toutes (fiches + exercices).
% Les gardes \ifdefined evitent tout conflit avec une redeclaration locale restee dans un
% fichier : ne pas retirer les gardes, ne pas redeclarer ces unites dans les fichiers.
\DeclareSIUnit{\jour}{j}
\ifdefined\var\else\DeclareSIUnit{\var}{var}\fi   % puissance reactive
\ifdefined\vad\else\DeclareSIUnit{\vad}{VAD}\fi   % puissance deformante
\ifdefined\tr\else\DeclareSIUnit{\tr}{tr}\fi      % tour (tr/min)
\usepackage{array, booktabs, multirow, tabularx}
\usepackage{xcolor}
\usepackage[most, breakable]{tcolorbox}
\usepackage{multicol}
\usepackage{enumitem}
\usepackage{mathtools}
\usepackage{lmodern}
% Refonte visuelle 2026 (direction "moderne coloré") : Lato pour le texte,
% maths sans-serif assorties via newtxsf (charge apres amsmath/amssymb).
\usepackage[default]{lato}
\usepackage{newtxsf}
\usepackage{microtype}
\usepackage{fancyhdr}
\usepackage{tikz}
\usetikzlibrary{positioning,arrows.meta}
\usepackage{circuitikz}
\usepackage{lastpage}

% ── Couleurs par domaine (palette modernisée 2026) ────────────────────
\definecolor{colElec}{RGB}{29, 78, 216}        % Électrotechnique — bleu vif
\definecolor{colElecLight}{RGB}{219, 234, 254} % fond bleu clair
\definecolor{colAuto}{RGB}{21, 128, 61}        % Automatique — vert
\definecolor{colMeca}{RGB}{194, 65, 12}        % Mécanique — orange
\definecolor{colTherm}{RGB}{185, 28, 28}       % Thermique — rouge
\definecolor{colQT}{RGB}{232, 244, 233}        % (hérité, plus utilisé par questiontype)
\definecolor{colJury}{RGB}{255, 248, 225}      % fond jury (crème)
\definecolor{colJuryFrame}{RGB}{202, 138, 4}   % bordure/bandeau jury (ambre)
\definecolor{colGrayLight}{RGB}{245, 245, 247} % fond gris neutre
\definecolor{colSeparator}{RGB}{180, 180, 180} % gris séparateur
\definecolor{colExo}{RGB}{8, 145, 178}          % Banque d'exercices — cyan (distinct des fiches)
% ── Encre et gris de la refonte ───────────────────────────────────────
\definecolor{colInk}{RGB}{23, 32, 42}          % texte courant (quasi-noir)
\definecolor{colSlate}{RGB}{83, 98, 118}       % gris ardoise (méta, en-têtes)
\definecolor{colRule}{RGB}{203, 213, 225}      % filets clairs
\definecolor{colJuryBg}{RGB}{255, 248, 225}    % fond crème jury
\definecolor{colPresence}{RGB}{83, 98, 118}    % bloc 5 (présence dans les sujets)
\AtBeginDocument{\color{colInk}}

% ── Style encadré exercice + corrigé (banque d'exercices) ─────────────
\tcbset{
  exoBox/.style={
    breakable,
    enhanced,
    colback=colExo!8!white,
    colframe=colExo,
    colbacktitle=colExo,
    coltitle=white,
    fonttitle=\bfseries\small,
    left=4pt, right=4pt, top=3pt, bottom=3pt,
    before skip=5pt, after skip=5pt,
  }
}

% ── Bascule enonce / corrige (banque d'exercices) ─────────────────────
% \CAPETCORRIGE est defini en ligne de commande (pdflatex "\def\CAPETCORRIGE{1}\input{...}")
% pour compiler la variante corriges d'un meme fichier source ; par defaut (non defini),
% la compilation produit la variante enonces. \CAPETSHOW{enonce}{corrige} choisit la branche.
%
% TROIS MODES (refonte 2026-09-20, manuel de revision) :
%   - aucun drapeau   -> imprime l'enonce            (variante enonces, INCHANGE)
%   - \CAPETCORRIGE   -> imprime le corrige          (variante corriges, INCHANGE)
%   - \CAPETMANUEL    -> imprime l'enonce ET met le corrige de cote pour l'annexe
%
% Le mode manuel stocke chaque corrige dans une macro (\csgdef), sans ecriture
% fichier ni re-tokenisation : le contenu (circuitikz, accents, %) traverse intact.
% L'annexe les rejoue ensuite dans l'ordre via \CAPETANNEXE.
\usepackage{etoolbox}

\newcounter{CAPETexo}          % numero continu d'exercice dans le manuel
\setcounter{CAPETexo}{0}
\newcounter{CAPETannexe}       % compteur de relecture en annexe

% Titre courant de la notion en cours (pose par \CAPETNOTION), sert d'intertitre annexe
\newcommand{\CAPETnotioncode}{}
\newcommand{\CAPETnotiontitre}{}

% \CAPETNOTION{code}{titre} : ouvre une notion dans le manuel (memorise le libelle,
% et marque le rang du premier exercice de la notion pour l'intertitre d'annexe).
\newcommand{\CAPETNOTION}[2]{%
  \renewcommand{\CAPETnotioncode}{#1}%
  \renewcommand{\CAPETnotiontitre}{#2}%
  \csgdef{CAPETnotionstart\the\numexpr\value{CAPETexo}+1\relax}{#1 -- #2}%
}

% ── \degree robuste ───────────────────────────────────────────────────
% Certains fichiers d'exercices utilisent \degree. Ni le preambule ni siunitx
% ne le definissent, et \providecommand local ne suffit pas : sous hyperref,
% \degree existe deja comme commande factice (\ERROR), donc \providecommand
% s'abstient et l'appel echoue a la compilation du manuel (erreur invisible
% en compilation isolee). On le (re)definit donc inconditionnellement ici,
% apres le chargement des packages, pour tout le corpus.
\AtBeginDocument{\providecommand{\degree}{}\renewcommand{\degree}{\ensuremath{^\circ}}}

% Libelle decoratif : les fichiers d'exercices utilisent aussi \CAPETSHOW dans
% leur \fancyhead et leur bloc de titre, pour afficher "(enonces)" ou "(corriges)".
% Ces appels-la ne sont PAS des exercices : en mode manuel ils ne doivent ni
% consommer un numero, ni alimenter l'annexe. On les reconnait a leur premier
% argument, strictement egal au libelle ci-dessous.
\newcommand{\CAPETLIBELLEENONCES}{(\'enonc\'es)}

\ifdefined\CAPETMANUEL
  % --- Mode manuel : enonce en place, corrige differe vers l'annexe ---
  \newcommand{\CAPETSHOW}[2]{%
    \def\CAPETargun{#1}%
    \ifx\CAPETargun\CAPETLIBELLEENONCES
      #1%            appel decoratif : rendu tel quel, hors numerotation
    \else
      \CAPETSHOWEXO{#1}{#2}%
    \fi
  }
  \newcommand{\CAPETSHOWEXO}[2]{%
    \stepcounter{CAPETexo}%
    \csgdef{CAPETcorrige\theCAPETexo}{#2}%
    \edef\CAPETtmpnum{\theCAPETexo}%
    #1%
    \par\vspace{2pt}%
    {\footnotesize\itshape\color{colExo}%
     \hyperref[capet:cor:\CAPETtmpnum]{Corrige de l'exercice \CAPETtmpnum\ $\rightarrow$ p.~\pageref{capet:cor:\CAPETtmpnum}}%
     \label{capet:eno:\CAPETtmpnum}}%
  }
  % \CAPETANNEXE : rejoue tous les corriges captures, dans l'ordre, avec intertitres
  \newcommand{\CAPETANNEXE}{%
    \setcounter{CAPETannexe}{0}%
    \loop\ifnum\value{CAPETannexe}<\value{CAPETexo}%
      \stepcounter{CAPETannexe}%
      \edef\CAPETcurnum{\theCAPETannexe}%
      \ifcsname CAPETnotionstart\CAPETcurnum\endcsname
        \par\medskip
        {\normalsize\bfseries\color{colExo}\csuse{CAPETnotionstart\CAPETcurnum}}%
        \par\nobreak\smallskip
      \fi
      \begin{tcolorbox}[exoBox, breakable,
        title={Corrige de l'exercice \CAPETcurnum\
               \hfill\normalfont\footnotesize
               \hyperref[capet:eno:\CAPETcurnum]{enonce p.~\pageref{capet:eno:\CAPETcurnum}}}]%
      \label{capet:cor:\CAPETcurnum}%
      \csuse{CAPETcorrige\CAPETcurnum}%
      \end{tcolorbox}%
    \repeat
  }
\else
  % --- Modes historiques, comportement strictement inchange ---
  \newcommand{\CAPETSHOW}[2]{\ifdefined\CAPETCORRIGE #2\else #1\fi}
  \newcommand{\CAPETANNEXE}{}
\fi

% ── Difficulté 1 à 3 (étoiles pleines, amssymb) ───────────────────────
\newcommand{\difficulte}[1]{%
  \ifcase#1\relax
  \or $\bigstar$\;niveau 1/3%
  \or $\bigstar\bigstar$\;niveau 2/3%
  \or $\bigstar\bigstar\bigstar$\;niveau 3/3%
  \fi
}

% ── Style encadré extrait de document technique (banque d'exercices) ──
\tcbset{
  docTech/.style={
    breakable,
    enhanced,
    colback=colGrayLight,
    colframe=colSeparator,
    fonttitle=\bfseries\small\itshape,
    coltitle=black,
    colbacktitle=colGrayLight,
    left=4pt, right=4pt, top=3pt, bottom=3pt,
    before skip=4pt, after skip=4pt,
  }
}

% ── Style tcolorbox de section principale ────────────────────────────
% Refonte 2026 : bandeau titre coloré arrondi + fond très pâle, sans cadre.
\tcbset{
  ficheSection/.style={
    breakable, enhanced jigsaw,
    frame hidden, boxrule=0pt, arc=1.8mm,
    colback=#1!4!white,
    colbacktitle=#1,
    coltitle=white,
    fonttitle=\sffamily\bfseries,
    left=7pt, right=7pt, top=5pt, bottom=5pt,
    toptitle=3pt, bottomtitle=3pt,
    before skip=10pt, after skip=8pt,
  }
}

% ── Style encadré question type ──────────────────────────────────────
% Carte blanche, bandeau sombre, liseré très clair.
\tcbset{
  questiontype/.style={
    breakable, enhanced jigsaw,
    colback=white,
    frame hidden, boxrule=0pt, arc=1.8mm,
    colbacktitle=colElec!78!black,
    coltitle=white,
    fonttitle=\sffamily\bfseries\small,
    left=7pt, right=7pt, top=4pt, bottom=4pt,
    toptitle=2.5pt, bottomtitle=2.5pt,
    before skip=8pt, after skip=8pt,
    borderline={0.6pt}{0pt}{colElec!12!white},
  }
}

% ── Style encadré jury ───────────────────────────────────────────────
% Bandeau ambre, fond crème, sans cadre.
\tcbset{
  jury/.style={
    breakable, enhanced jigsaw,
    colback=colJuryBg,
    frame hidden, boxrule=0pt, arc=1.8mm,
    colbacktitle=colJuryFrame,
    coltitle=white,
    fonttitle=\sffamily\bfseries,
    left=7pt, right=7pt, top=4pt, bottom=4pt,
    toptitle=3pt, bottomtitle=3pt,
    before skip=10pt, after skip=8pt,
  }
}

% ── En-tête / pied de page communs (refonte 2026) ─────────────────────
% #1 code fiche, #2 titre court, #3 couleur domaine
\newcommand{\ficheHeaders}[3]{%
  \fancyhead[L]{\small\sffamily{\bfseries\textcolor{#3}{#1}}\hspace{0.7em}\textcolor{colSlate}{#2}}%
  \fancyhead[R]{\small\sffamily\textcolor{colSlate}{CAPET SII IE\, --- \,Session 2026}}%
  \fancyfoot[C]{\small\sffamily\textcolor{colSlate}{\thepage\,/\,\pageref{LastPage}}}%
  \fancyhead[C]{}\fancyfoot[L]{}\fancyfoot[R]{}%
}

% ── Bloc en-tête de fiche (BLOC 0, refonte 2026) ──────────────────────
% #1 code  #2 titre  #3 couleur domaine  #4 domaine  #5 ligne présence  #6 sous-thèmes
\newcommand{\ficheEntete}[6]{%
  \begin{tcolorbox}[enhanced, breakable,
    frame hidden, boxrule=0pt, arc=2.5mm,
    interior style={left color=#3, right color=#3!75!black},
    left=9pt, right=9pt, top=7pt, bottom=7pt,
    before skip=0pt, after skip=6pt]
    {\color{white}\sffamily\bfseries\LARGE #1\quad #2}\\[4pt]
    {\color{#3!22!white}\sffamily\small \MakeUppercase{#4}\quad\textbar\quad #5}\\[1pt]
    {\color{#3!35!white}\sffamily\footnotesize Sous-thèmes : #6}
  \end{tcolorbox}%
}

% ── Mise en page ─────────────────────────────────────────────────────
\pagestyle{fancy}
\fancyhf{}
\renewcommand{\headrulewidth}{0.5pt}
\renewcommand{\footrulewidth}{0pt}
\AtBeginDocument{%
  \renewcommand{\headrule}{{\color{colRule}\hrule width\headwidth height 0.5pt}}%
}

\setlist[itemize]{leftmargin=1.2em, itemsep=1pt, topsep=2pt, parsep=0pt}
\setlist[enumerate]{leftmargin=1.2em, itemsep=3pt, topsep=2pt, parsep=0pt}

\setlength{\parskip}{2pt}
\setlength{\columnsep}{10pt}

\def\CAPETBEGINFILE{\begin{document}}
\def\CAPETENDFILE{\end{document}}
\else
\def\CAPETBEGINFILE{}
\def\CAPETENDFILE{}
\fi

% Accolades du schema de troncature binaire (exercice 2) : la bibliotheque
% decorations.pathreplacing n'est pas charge par ../fiches/preamble.tex.
% Symbole triangulaire de l'ALI (exercices 4 et 6) : shapes.geometric.
\usetikzlibrary{decorations.pathreplacing, shapes.geometric}

\fancyhead[L]{\textbf{\textcolor{colExo}{T22}} - Exercices - CAN/CNA et numerisation}
\fancyhead[R]{\small\textit{Banque d'exercices \CAPETSHOW{(\'enonc\'es)}{(corrig\'es)}}}
\fancyfoot[C]{\small\thepage\,/\,\pageref{LastPage}}

\CAPETBEGINFILE

\begin{tcolorbox}[
  enhanced, breakable,
  colback=colExo,
  colframe=colExo,
  left=6pt, right=6pt, top=5pt, bottom=5pt,
  before skip=0pt, after skip=4pt
]
  {\color{white}\bfseries\large T22 \quad CAN/CNA, num\'erisation et \'echantillonnage}\\[2pt]
  {\color{white}\small Banque d'exercices -- 6 exercices \CAPETSHOW{(\'enonc\'es)}{(corrig\'es)}}
\end{tcolorbox}

% =====================================================================
% Exercice 1 -- 2022, questions 23 et 24
% =====================================================================
\begin{tcolorbox}[exoBox,
  title={Exercice 1 -- \difficulte{1} \hfill\normalfont\footnotesize 2022, \'epreuve \'ecrite disciplinaire (option IE), questions 23 et 24}
]
\CAPETSHOW{
Le module de t\'el\'etransmission TRT~200 d'un poste de purge de r\'eseau de gaz surveille en permanence son alimentation \'electrique. Trois grandeurs analogiques sont appliqu\'ees sur les entr\'ees \textbf{ADC0}, \textbf{ADC1} et \textbf{ADC2} du convertisseur analogique-num\'erique int\'egr\'e au microcontr\^oleur PIC.

Le sujet pr\'ecise que ce CAN est un convertisseur \textbf{12 bits} et que sa \textbf{tension de r\'ef\'erence vaut $V_{REF+} = \SI{3,5}{\volt}$} (avec $V_{REF-} = \SI{0}{\volt}$). L'entr\'ee \textbf{ADC2} surveille la tension $U_{Peltier}$ d\'elivr\'ee par le module thermo\'electrique~; elle lui est reli\'ee par un \textbf{pont diviseur de rapport $1/2$} form\'e de deux r\'esistances \'egales $R_4$ et $R_5$.

Le diagramme des exigences impose que la tension du module Peltier ne descende \textbf{jamais en dessous de \SI{0,4}{\volt}} : en de\c{c}\`a, une alarme doit \^etre t\'el\'etransmise.

\begin{center}
\begin{circuitikz}[european, scale=0.9, every node/.style={scale=0.9}]
  % source Peltier
  \draw (0,0) to[V, l=$U_{Peltier}$, invert] (0,3);
  \draw (0,3) -- (2,3) node[circ]{};
  % pont diviseur
  \draw (2,3) to[R, l=$R_5$] (2,1.7);
  \draw (2,1.7) node[circ]{} to[R, l=$R_4$] (2,0);
  \draw (0,0) -- (2,0);
  % point milieu vers CAN
  \draw (2,1.7) -- (4.2,1.7) node[midway,above,font=\scriptsize]{ADC2};
  % bloc CAN
  \node[draw, minimum width=2.3cm, minimum height=1.5cm, align=center, font=\scriptsize] (can) at (5.4,1.7) {CAN\\12 bits\\$V_{REF+}{=}\SI{3,5}{\volt}$};
  \draw[-latex] (can.east) -- ++(1.3,0) node[right,font=\scriptsize]{$N$};
  \draw (2,0) -- (5.4,0) -- (can.south);
  \node[ground] at (3.6,0) {};
\end{circuitikz}
\end{center}

\begin{enumerate}[leftmargin=1.5em]
  \item Pr\'eciser la r\'esolution de ce convertisseur, puis \textbf{calculer son quantum} $q$ en \si{\milli\volt\per\text{LSB}}. On adoptera la convention retenue par le corrig\'e officiel du jury, $q = (V_{REF+} - V_{REF-})/(2^n - 1)$, et l'on justifiera qu'\`a cette r\'esolution le choix de la convention est sans cons\'equence pratique.
  \item En d\'eduire la valeur \textbf{d\'ecimale} $N$ d\'elivr\'ee en sortie du CAN lorsque la tension aux bornes du module Peltier vaut $U_{Peltier} = \SI{2}{\volt}$.
  \item Lors d'une num\'erisation, la sortie du CAN vaut $\mathtt{0xAF}$. \textbf{Justifier} l'alarme t\'el\'etransmise par le module.
\end{enumerate}

}{
\textbf{Corrig\'e.}

\begin{enumerate}[leftmargin=1.5em]
  \item \textbf{R\'esolution.} La r\'esolution est le \textbf{nombre de bits} du mot de sortie, ici $n = \textbf{12 bits}$, soit $2^{12} = 4\,096$ codes distincts (de $0$ \`a $4\,095$).

  \emph{Attention} : la r\'esolution est un \textbf{nombre de bits}, pas une tension -- c'est la confusion la plus sanctionn\'ee sur ce type de question. La tension associ\'ee \`a un pas de code est le \textbf{quantum} :
  \[ q = \frac{V_{REF+} - V_{REF-}}{2^n - 1} = \frac{3{,}5 - 0}{4\,095} = \SI{8,547e-4}{\volt} \approx \boldsymbol{\SI{0,855}{\milli\volt\per\text{LSB}}} \]

  \emph{Choix de la convention.} Avec l'autre convention usuelle, $q = V_{REF}/2^n = 3{,}5/4\,096 = \SI{0,8545}{\milli\volt}$. L'\'ecart relatif entre les deux vaut $1/(2^n-1) = 1/4\,095 \approx \SI{0,024}{\percent}$ : totalement n\'egligeable devant les erreurs de lin\'earit\'e et d'offset d'un CAN r\'eel. Le choix de la convention n'a de cons\'equence pratique que pour les tr\`es faibles r\'esolutions ($n \leq 4$). Ce qui compte en revanche, c'est de \textbf{l'annoncer} avant de calculer.
  \item Le pont diviseur \'etant de rapport $1/2$, la tension r\'eellement pr\'esent\'ee \`a l'entr\'ee du CAN vaut :
  \[ V_{ADC2} = \frac{U_{Peltier}}{2} = \frac{2}{2} = \SI{1}{\volt} \]
  \[ N = \frac{V_{ADC2}}{q} = \frac{1}{8{,}547\times 10^{-4}} = \boldsymbol{1\,170} \]
  \emph{V\'erification de coh\'erence} : $1\,170 \times 0{,}8547\times 10^{-3} = \SI{1,000}{\volt}$ -- on retrouve bien la tension de d\'epart. Le code $1\,170$ occupe environ \SI{29}{\percent} de la dynamique du convertisseur.
  \item Conversion du code hexad\'ecimal en d\'ecimal :
  \[ \mathtt{0xAF} = \mathrm{A}_{16}\times 16 + \mathrm{F}_{16} = 10\times 16 + 15 = 175_{10} \qquad (= 1010\,1111_2) \]
  On \textbf{remonte} la cha\^ine, du code vers la grandeur physique :
  \[ V_{ADC2} = N\cdot q = 175\times 8{,}547\times 10^{-4} = \SI{0,1496}{\volt} \approx \SI{0,15}{\volt} \]
  \[ U_{Peltier} = 2\times V_{ADC2} = 2\times 0{,}1496 = \boldsymbol{\SI{0,299}{\volt}} \approx \SI{0,3}{\volt} \]

  \smallskip
  \textbf{Conclusion.} La tension reconstitu\'ee aux bornes du module Peltier, $U_{Peltier} \approx \SI{0,30}{\volt}$, est \textbf{inf\'erieure au seuil de \SI{0,4}{\volt}} impos\'e par le diagramme des exigences : le module Peltier ne fournit plus une tension suffisante pour alimenter la carte \'electronique du TRT~200. L'alarme t\'el\'etransmise est donc \textbf{justifi\'ee}, et l'installation doit basculer sur les accumulateurs de secours. Le code d'alarme correspondant est $N_{seuil} = 0{,}4/2/q \approx 234$, soit $\mathtt{0xEA}$ : tout code inf\'erieur \`a $234$ doit d\'eclencher l'alarme.
\end{enumerate}
}
\end{tcolorbox}

% =====================================================================
% Exercice 2 -- 2023, questions D2 et D3
% =====================================================================
\begin{tcolorbox}[exoBox,
  title={Exercice 2 -- \difficulte{1} \hfill\normalfont\footnotesize 2023, \'epreuve \'ecrite disciplinaire (option IE), questions D2 et D3}
]
\CAPETSHOW{
La perche porte-outils d'un enjambeur viticole autonome est d\'eplac\'ee par deux v\'erins \'electriques. La position de la tige de v\'erin est mesur\'ee par un \textbf{codeur incr\'emental} dont la r\'esolution est de \textbf{\SI{0,15}{\milli\meter} par impulsion}.

\begin{itemize}[leftmargin=1.5em, itemsep=1pt]
  \item l'excursion maximale de la tige du v\'erin est de \textbf{\SI{680}{\milli\meter}} ;
  \item lorsque le v\'erin est compl\`etement rentr\'e, la valeur d\'ecimale du compteur d'impulsions est \textbf{0} ;
  \item la \textbf{pr\'ecision requise} pour contr\^oler la position de la perche est de \textbf{\SI{5}{\milli\meter}}.
\end{itemize}

La carte d'asservissement, pour all\'eger son calcul, ne conserve que les \textbf{8 bits de poids fort} du mot de position.

\begin{center}
\begin{tikzpicture}[>=Stealth, scale=0.92, font=\footnotesize]
  % mot de 13 bits
  \foreach \i in {0,...,12} {
    \draw (\i*0.62,0) rectangle ++(0.62,0.62);
  }
  % 8 MSB surlignes
  \fill[colExo!25] (0,0) rectangle (8*0.62,0.62);
  \foreach \i in {0,...,12} {
    \draw (\i*0.62,0) rectangle ++(0.62,0.62);
  }
  \node[font=\scriptsize] at (0.31,0.31) {$b_{12}$};
  \node[font=\scriptsize] at (0.93,0.31) {$b_{11}$};
  \node[font=\scriptsize] at (1.55,0.31) {$\cdots$};
  \node[font=\scriptsize] at (4.65,0.31) {$b_{5}$};
  \node[font=\scriptsize] at (5.27,0.31) {$b_{4}$};
  \node[font=\scriptsize] at (5.89,0.31) {$\cdots$};
  \node[font=\scriptsize] at (7.75,0.31) {$b_{0}$};
  % accolades
  \draw[decorate, decoration={brace, amplitude=4pt}] (0,0.72) -- (8*0.62,0.72)
    node[midway, above, font=\scriptsize, yshift=2pt]{8 bits de poids fort \emph{conserv\'es}};
  \draw[decorate, decoration={brace, mirror, amplitude=4pt}] (8*0.62,-0.1) -- (13*0.62,-0.1)
    node[midway, below, font=\scriptsize, yshift=-2pt]{5 bits \emph{supprim\'es}};
  \node[left, font=\scriptsize] at (-0.15,0.31) {MSB};
  \node[right, font=\scriptsize] at (13*0.62+0.15,0.31) {LSB};
\end{tikzpicture}
\end{center}

\begin{enumerate}[leftmargin=1.5em]
  \item \textbf{D\'eterminer} le nombre de bits n\'ecessaires pour coder en binaire la position de la tige du v\'erin sur toute son excursion. \textbf{Indiquer} le code de la position maximale, en binaire.
  \item \textbf{D\'eterminer} la pr\'ecision de mesure, en \si{\milli\meter}, obtenue en ne retenant que les 8 bits de poids fort. \textbf{Justifier} que ce choix est compatible avec la mise en \oe uvre de l'asservissement.
  \item \textbf{Donner} l'octet correspondant \`a la position maximale, en binaire et en hexad\'ecimal.
\end{enumerate}

}{
\textbf{Corrig\'e.}

\begin{enumerate}[leftmargin=1.5em]
  \item \textbf{Nombre d'impulsions \`a coder.} Le codeur d\'elivre une impulsion tous les \SI{0,15}{\milli\meter} ; sur l'excursion compl\`ete :
  \[ N_{imp} = \frac{680}{0{,}15} = 4\,533{,}3 \quad\Longrightarrow\quad \text{code maximal} = 4\,533 \]
  \textbf{Nombre de bits.} Il faut pouvoir coder toutes les valeurs de $0$ \`a $4\,533$, soit $4\,534$ valeurs distinctes :
  \[ 2^{12} = 4\,096 < 4\,534 \leq 8\,192 = 2^{13} \quad\Longrightarrow\quad \boldsymbol{n = 13\ \text{bits}} \]
  \emph{Pi\`ege classique} : $\log_2(4\,533) = 12{,}15$, ce qui incite \`a r\'epondre « 12 bits » par troncature. Il faut \textbf{arrondir \`a l'entier sup\'erieur} : avec 12 bits on ne coderait que jusqu'\`a $4\,095$, soit \SI{614}{\milli\meter}, et les \SI{66}{\milli\meter} de fin de course seraient inaccessibles.

  \textbf{Code de la position maximale} sur 13 bits :
  \[ 4\,533_{10} = \boldsymbol{1\,0001\,1011\,0101_2} \]
  \emph{V\'erification} : $4\,096 + 256 + 128 + 32 + 16 + 4 + 1 = 4\,533$. \checkmark
  \item Ne conserver que les \textbf{8 bits de poids fort} d'un mot de \textbf{13 bits}, c'est \textbf{supprimer les $13 - 8 = 5$ bits de poids faible}, donc diviser le code par $2^5 = 32$. Chaque incr\'ement du mot tronqu\'e correspond alors \`a 32 impulsions du codeur :
  \[ \delta = 32 \times 0{,}15 = \boldsymbol{\SI{4,8}{\milli\meter}} \]
  \textbf{Justification.} La pr\'ecision requise \'etant de \SI{5}{\milli\meter} :
  \[ \delta = \SI{4,8}{\milli\meter} \;<\; \SI{5}{\milli\meter} \quad\checkmark \]
  Le choix est donc \textbf{compatible}, mais avec une marge tr\`es faible (\SI{4}{\percent}). L'int\'er\^et de la troncature est de ramener la mesure de position sur un \textbf{unique octet}, directement manipulable par le microcontr\^oleur 8 bits de la carte de contr\^ole : un seul registre \`a lire, aucune op\'eration multi-octets dans la boucle d'asservissement, donc un temps de cycle plus court. Conserver les 13 bits n'apporterait qu'une r\'esolution inutile au regard de l'exigence.

  \emph{Remarque} : c'est bien le \textbf{poids fort} qu'il faut conserver. Garder les 8 bits de poids \emph{faible} reviendrait \`a perdre toute information sur la position absolue (le mot repasserait par z\'ero tous les \SI{38,4}{\milli\meter}).
  \item On d\'ecale le mot de 13 bits de 5 rangs vers la droite :
  \[ \underbrace{1\,0001\,101}_{8\ \text{MSB}}\underbrace{1\,0101}_{\text{supprim\'es}} \quad\Longrightarrow\quad \text{octet} = 1000\,1101_2 \]
  \[ 1000\,1101_2 = 128 + 8 + 4 + 1 = 141_{10} = \boldsymbol{\mathtt{8D}_{16}} \]
  \emph{V\'erification de coh\'erence} : $141 \times 32 \times 0{,}15 = \SI{676,8}{\milli\meter}$, \`a comparer \`a l'excursion r\'eelle de \SI{680}{\milli\meter}. L'\'ecart de \SI{3,2}{\milli\meter} est bien inf\'erieur au pas de \SI{4,8}{\milli\meter} : c'est simplement l'effet de la troncature, qui arrondit toujours \textbf{par d\'efaut}.

  \smallskip
  \textbf{Conclusion.} Le codage sur \textbf{13 bits} couvre l'int\'egralit\'e de la course de \SI{680}{\milli\meter}, et la troncature aux \textbf{8 bits de poids fort} ram\`ene la position sur un octet unique ($\mathtt{8D}_{16}$ en but\'ee haute) avec une pr\'ecision de \SI{4,8}{\milli\meter}, \textbf{conforme \`a l'exigence de \SI{5}{\milli\meter}} du cahier des charges. La marge \'etant faible, cette troncature ne pourra pas \^etre reconduite telle quelle si l'exigence de pr\'ecision \'etait resserr\'ee : il faudrait alors conserver 9 bits (pas de \SI{2,4}{\milli\meter}).
\end{enumerate}
}
\end{tcolorbox}

% =====================================================================
% Exercice 3 -- docTech -- 2026, questions B10 et B12
% =====================================================================
\begin{tcolorbox}[exoBox,
  title={Exercice 3 -- \difficulte{2} \hfill\normalfont\footnotesize 2026, \'epreuve 2 d'admissibilit\'e (BAC + 3, option IE), questions B10 et B12}
]
\CAPETSHOW{
La temp\'erature de fonctionnement d'une pile \`a combustible embarqu\'ee dans un autobus est mesur\'ee par une cha\^ine d'acquisition compl\`ete : capteur Jumo, conditionneur, amplificateur, puis convertisseur analogique-num\'erique.

Le \textbf{cahier des charges de la mesure} est le suivant :

\begin{center}
\begin{tabular}{@{}ll@{}}
\toprule
\'Etendue de mesure & $-5\,^\circ$C \`a $+90\,^\circ$C \\
Erreur de mesure & $< 1\,^\circ$C \\
\bottomrule
\end{tabular}
\end{center}

L'ensemble \{amplificateur + conditionneur\} d\'elivre une tension $v_A$ \textbf{directement proportionnelle \`a la temp\'erature}, $v_A = c\cdot T$, \'evoluant sur une \textbf{plage de tension de \SI{3}{\volt}} pour couvrir toute l'\'etendue de mesure. Le CAN choisi par le constructeur est le \textbf{Texas Instruments ADS7029-Q1}, aliment\'e sous $AVDD = \SI{3}{\volt}$.

\begin{tcolorbox}[docTech, title={Document technique -- Extrait de la notice constructeur ADS7029-Q1 (Texas Instruments, SBAS811)}]
\footnotesize
\textbf{ADS7029-Q1 -- Small-Size, Low-Power, 8-Bit, 2-MSPS, SAR ADC}

\smallskip
\textbf{Caract\'eristiques \'electriques} ($T_A = -40$ \`a $+125\,^\circ$C, $AVDD = \SI{3}{\volt}$, $f_{SAMPLE} = 2$~MSPS) :

\begin{center}
\begin{tabular}{@{}llcccl@{}}
\toprule
& Param\`etre & Min & Typ & Max & Unit\'e \\
\midrule
\multicolumn{6}{@{}l}{\textit{Entr\'ee analogique}} \\
& \emph{Full-scale input voltage span} & $0$ & & $AVDD$ & V \\
& \emph{Reference voltage} $V_{REF} = AVDD$ & $2{,}35$ & & $3{,}6$ & V \\
\midrule
\multicolumn{6}{@{}l}{\textit{Performances syst\`eme}} \\
& \emph{Resolution} & & 8 & & bits \\
& \emph{No missing codes} (NMC) & 8 & & & bits \\
INL & \emph{Integral nonlinearity} & $-0{,}5$ & $\pm 0{,}25$ & $+0{,}5$ & LSB \\
DNL & \emph{Differential nonlinearity} & $-0{,}4$ & $\pm 0{,}2$ & $+0{,}4$ & LSB \\
$E_O$ & \emph{Offset error} & & $\pm 0{,}5$ & & LSB \\
$E_G$ & \emph{Gain error} & & $\pm 0{,}2$ & & \%FS \\
\midrule
\multicolumn{6}{@{}l}{\textit{Dynamique d'\'echantillonnage}} \\
$t_{ACQ}$ & \emph{Acquisition time} & 120 & & & ns \\
$t_{CYCLE}$ & \emph{Cycle time} & & $0{,}5$ & & \si{\micro\second} \\
$f_{THROUGHPUT}$ & \emph{Maximum throughput rate} & & 2 & & MSPS \\
\bottomrule
\end{tabular}
\end{center}

\smallskip
\textbf{\S8.3.3 -- ADC Transfer Function.} \emph{The device output is in straight binary format. The device resolution for a single-ended input can be computed by Equation 2 :}
\[ 1\ \mathrm{LSB} = V_{REF} / 2^N \qquad \text{avec } V_{REF} = AVDD \text{ et } N = 8 \]

\smallskip
\textbf{Rappel du document DTB5 du sujet.} \emph{La courbe de transfert th\'eorique donnant le mot binaire de sortie $N_{(2)}$ en fonction de la tension d'entr\'ee $V_E$ est une courbe en escalier, dans le cas d'une quantification centr\'ee. On peut lui associer l'erreur de conversion ou erreur de quantification, en dents de scie. Celle-ci est comprise entre $-q/2$ et $+q/2$.}
\end{tcolorbox}

\begin{enumerate}[leftmargin=1.5em]
  \item \textbf{Relever} dans le document technique les trois caract\'eristiques du CAN n\'ecessaires \`a la suite de l'\'etude : la \textbf{grandeur pleine \'echelle}, la \textbf{r\'esolution} et la \textbf{formule de d\'efinition du quantum} retenue par le constructeur. En d\'eduire la valeur num\'erique du \textbf{quantum} $q$.
  \item La litt\'erature emploie aussi la convention $q = V_{REF}/(2^n - 1)$. \textbf{Calculer} le quantum avec cette seconde convention et \textbf{quantifier} l'\'ecart relatif entre les deux. Conclure sur la convention \`a retenir ici.
  \item \textbf{D\'eterminer} l'erreur de quantification maximale du CAN, en volt.
  \item Sachant que la plage de \SI{3}{\volt} couvre exactement l'\'etendue de mesure $[-5\,;\,90]\,^\circ$C, \textbf{d\'eterminer} le coefficient $c$ de la cha\^ine, puis \textbf{en d\'eduire} l'erreur de mesure de la temp\'erature $T$ due \`a la seule quantification. \textbf{Conclure} vis-\`a-vis du cahier des charges.
  \item Le calcul pr\'ec\'edent ne retient que l'erreur de quantification. \`A l'aide du document technique, \textbf{citer} deux autres sources d'erreur du convertisseur et \textbf{estimer} l'erreur totale de la cha\^ine dans le pire cas.
\end{enumerate}

}{
\textbf{Corrig\'e.}

\textbf{\'Etape 1 -- Extraction des donn\'ees du document technique.}

\begin{enumerate}[leftmargin=1.5em]
  \item Lecture directe de la notice :
  \begin{itemize}[leftmargin=1.5em, itemsep=1pt]
    \item \textbf{Grandeur pleine \'echelle} : la ligne \emph{Full-scale input voltage span} donne $[0\,;\,AVDD]$. Le CAN \'etant aliment\'e sous $AVDD = \SI{3}{\volt}$, la pleine \'echelle vaut $PE = \boldsymbol{\SI{3}{\volt}}$ -- ce qui co\"incide exactement avec la plage de \SI{3}{\volt} impos\'ee \`a l'ensemble \{amplificateur + conditionneur\} : la cha\^ine exploite bien \textbf{toute} la dynamique du convertisseur.
    \item \textbf{R\'esolution} : $\boldsymbol{n = 8\ \text{bits}}$ (titre de la notice et ligne \emph{Resolution}), soit $2^8 = 256$ codes. La ligne \emph{No missing codes} = 8 bits garantit qu'aucun de ces 256 codes n'est manquant.
    \item \textbf{Convention du constructeur} : le \S8.3.3 impose $1\ \mathrm{LSB} = V_{REF}/2^N$, donc la convention en $\boldsymbol{2^n}$.
  \end{itemize}
  \[ q = \frac{V_{REF}}{2^n} = \frac{3}{256} = \SI{1,1719e-2}{\volt} \approx \boldsymbol{\SI{11,72}{\milli\volt}} \]
  \item Avec la seconde convention :
  \[ q' = \frac{V_{REF}}{2^n - 1} = \frac{3}{255} \approx \SI{11,76}{\milli\volt} \]
  \[ \frac{q' - q}{q} = \frac{1}{2^n - 1} = \frac{1}{255} \approx \SI{0,39}{\percent} \]
  L'\'ecart de \SI{0,39}{\percent} est \textbf{n\'egligeable devant les erreurs propres du composant} relev\'ees dans la notice (INL jusqu'\`a $\pm 0{,}5$~LSB, soit \SI{50}{\percent} d'un quantum, et erreur de gain $\pm 0{,}2$~\%FS). On retient donc ici la convention \textbf{du constructeur}, $q = V_{REF}/2^n$ : lorsqu'un document technique impose une d\'efinition, c'est elle qui fait foi. Dans tous les cas, la convention retenue doit \^etre \textbf{annonc\'ee avant le calcul}.

  \emph{Origine de l'\'ecart} : $2^n - 1$ compte le nombre d'\emph{intervalles} entre les $2^n$ codes (la derni\`ere transition a lieu \`a $V_{REF} - 1\,\mathrm{LSB}$), tandis que $2^n$ d\'ecoupe la pleine \'echelle en $2^n$ plages de largeur \'egale. Les deux lectures sont d\'efendables ; seule l'incoh\'erence entre l'annonce et le calcul est sanctionn\'ee.
  \item Le document DTB5 pr\'ecise que la quantification est \textbf{centr\'ee} : la caract\'eristique en escalier est d\'ecal\'ee d'un demi-quantum, si bien que l'erreur de conversion est born\'ee sym\'etriquement :
  \[ \varepsilon_{max} = \pm\frac{q}{2} = \pm\frac{11{,}72}{2} = \boldsymbol{\pm\SI{5,86}{\milli\volt}} \]

  \begin{center}
  \begin{tikzpicture}[scale=0.95, font=\footnotesize]
    % --- caracteristique en escalier ---
    \draw[->] (0,0) -- (7.4,0) node[right] {$V_E$};
    \draw[->] (0,0) -- (0,3.5) node[above] {$N$};
    % droite ideale
    \draw[thick, colExo!40] (0,0) -- (7,3.0);
    % escalier (quantification centree : premiere marche de largeur q/2)
    \draw[thick, colExo]
      (0,0) -- (0.5,0) -- (0.5,0.5) -- (1.5,0.5) -- (1.5,1.0) -- (2.5,1.0)
      -- (2.5,1.5) -- (3.5,1.5) -- (3.5,2.0) -- (4.5,2.0) -- (4.5,2.5)
      -- (5.5,2.5) -- (5.5,3.0) -- (6.5,3.0);
    % quantum
    \draw[<->] (6.8,2.5) -- (6.8,3.0);
    \node[right, font=\scriptsize] at (6.85,2.75) {$q$};
    \draw[dotted] (2.5,0) -- (2.5,1.5);
    \draw[dotted] (3.5,0) -- (3.5,2.0);
    \draw[<->] (2.5,-0.32) -- (3.5,-0.32) node[midway, below, font=\scriptsize]{$q$};
    \node[below, font=\scriptsize] at (0.25,-0.05) {$0$};
    % --- erreur en dents de scie ---
    \begin{scope}[yshift=-2.55cm]
      \draw[->] (0,0) -- (7.4,0) node[right] {$V_E$};
      \draw[->] (0,-0.75) -- (0,0.9) node[above, font=\scriptsize] {$\varepsilon$};
      \draw[dashed, colSlate] (0,0.5) -- (7,0.5) node[right, font=\scriptsize]{$+q/2$};
      \draw[dashed, colSlate] (0,-0.5) -- (7,-0.5) node[right, font=\scriptsize]{$-q/2$};
      \draw[thick, colTherm]
        (0,0) -- (0.5,0.5) -- (0.5,-0.5) -- (1.5,0.5) -- (1.5,-0.5)
        -- (2.5,0.5) -- (2.5,-0.5) -- (3.5,0.5) -- (3.5,-0.5)
        -- (4.5,0.5) -- (4.5,-0.5) -- (5.5,0.5) -- (5.5,-0.5) -- (6.5,0.5);
    \end{scope}
  \end{tikzpicture}
  \end{center}
  Caract\'eristique en escalier (haut) et erreur de quantification en dents de scie (bas), born\'ee \`a $\pm q/2$.
  \item La plage de \SI{3}{\volt} couvre l'\'etendue $[-5\,;\,90]\,^\circ$C, soit $\Delta T = 90 - (-5) = 95\,^\circ$C :
  \[ c = \frac{PE}{\Delta T} = \frac{3}{95} \approx \boldsymbol{\SI{0,0316}{\volt\per\degreeCelsius}} \]
  L'erreur en temp\'erature s'obtient en \textbf{divisant} l'erreur en tension par la sensibilit\'e de la cha\^ine (op\'eration inverse de la conversion temp\'erature $\rightarrow$ tension) :
  \[ \varepsilon_T = \frac{q/2}{c} = \frac{5{,}86\times 10^{-3}}{0{,}0316} \approx \boldsymbol{\pm\SI{0,19}{\degreeCelsius}} \]
  \emph{V\'erification par une autre voie} : l'\'etendue de $95\,^\circ$C est d\'ecoup\'ee en $256$ pas, soit $95/256 = 0{,}371\,^\circ$C par pas ; la moiti\'e vaut $0{,}186\,^\circ$C. \checkmark

  \smallskip
  \textbf{Conclusion.} L'erreur de mesure due \`a la quantification, $\pm\SI{0,19}{\degreeCelsius}$, est \textbf{cinq fois plus petite} que l'erreur de mesure maximale de $1\,^\circ$C impos\'ee par le cahier des charges : le \textbf{choix d'un CAN 8 bits est suffisant} et laisse une marge confortable pour les autres sources d'erreur de la cha\^ine (capteur, conditionneur, d\'erives). Monter en r\'esolution serait ici une d\'epense inutile.
  \item \textbf{Autres sources d'erreur relev\'ees dans la notice.}
  \begin{itemize}[leftmargin=1.5em, itemsep=1pt]
    \item l'\textbf{INL} (non-lin\'earit\'e int\'egrale), jusqu'\`a $\pm 0{,}5$~LSB au pire cas, soit $\pm\SI{5,86}{\milli\volt}$ : \'ecart entre la caract\'eristique r\'eelle et la droite id\'eale ;
    \item l'\textbf{erreur d'offset} $E_O = \pm 0{,}5$~LSB, soit encore $\pm\SI{5,86}{\milli\volt}$ ;
    \item (accessoirement l'\textbf{erreur de gain} $E_G = \pm 0{,}2$~\%FS $= \pm\SI{6}{\milli\volt}$ sur la pleine \'echelle).
  \end{itemize}
  Dans le \textbf{pire cas}, en sommant arithm\'etiquement quantification, INL et offset :
  \[ \varepsilon_{tot} \approx 5{,}86 + 5{,}86 + 5{,}86 = \SI{17,6}{\milli\volt} \quad\Longrightarrow\quad \varepsilon_{T,tot} = \frac{17{,}6\times 10^{-3}}{0{,}0316} \approx \SI{0,56}{\degreeCelsius} \]
  \textbf{Conclusion compl\'ementaire.} M\^eme en cumulant les trois d\'efauts au pire cas, l'erreur reste inf\'erieure \`a $1\,^\circ$C : le convertisseur \textbf{satisfait seul} l'exigence, et le budget d'erreur restant ($1 - 0{,}56 = 0{,}44\,^\circ$C) est disponible pour le capteur et le conditionneur. C'est ce raisonnement de \textbf{budget d'erreur} -- et non le seul calcul du quantum -- qui est attendu \`a ce niveau.

  \smallskip
  \textit{Point de vigilance jury (rapport de jury CAPET 2020, cinqui\`eme partie) : \guillemotleft\ Le jury d\'eplore que la majeure partie des candidats ne sache pas d\'eterminer la r\'esolution d'un capteur. \guillemotright\ Rappelons la distinction : la \textbf{r\'esolution} du CAN est un nombre de \textbf{bits} (8 ici) ; le \textbf{quantum} est la tension d'un pas (\SI{11,72}{\milli\volt}) ; la \textbf{r\'esolution de la cha\^ine} ramen\'ee \`a la grandeur physique est $q/c = \SI{0,37}{\degreeCelsius}$ ; la \textbf{pr\'ecision} enfin est la borne d'erreur totale, ici $\pm\SI{0,56}{\degreeCelsius}$. Les quatre notions ne se confondent pas.}
\end{enumerate}
}
\end{tcolorbox}

% =====================================================================
% Exercice 4 -- 2026, questions B13, B14, B15 et B16
% =====================================================================
\begin{tcolorbox}[exoBox,
  title={Exercice 4 -- \difficulte{2} \hfill\normalfont\footnotesize 2026, \'epreuve 2 d'admissibilit\'e (BAC + 3, option IE), questions B13 \`a B16}
]
\CAPETSHOW{
On s'int\'eresse d\'esormais \`a l'\textbf{asservissement num\'erique} de la temp\'erature de la m\^eme pile \`a combustible. La cha\^ine directe comporte un correcteur num\'erique, un \textbf{CNA}, puis le syst\`eme de refroidissement ; la cha\^ine de retour comporte le capteur, un \textbf{filtre anti-repliement}, un \textbf{\'echantillonneur-bloqueur (E/B)} et le \textbf{CAN}.

\begin{center}
\begin{tikzpicture}[>=Stealth, scale=0.9, every node/.style={scale=0.9}, font=\scriptsize]
  \node[draw, circle, minimum size=0.5cm, inner sep=0pt] (som) at (0.9,1.9) {};
  \node[font=\scriptsize] at (0.62,2.18) {$+$};
  \node[font=\scriptsize] at (0.62,1.55) {$-$};
  \draw[->] (-0.3,1.9) -- (som.west) node[midway,above,font=\scriptsize]{$T_{cons}$};
  \node[draw, minimum width=1.5cm, minimum height=0.8cm, align=center] (cor) at (2.5,1.9) {Correcteur\\num\'erique};
  \node[draw, minimum width=0.9cm, minimum height=0.8cm] (cna) at (4.4,1.9) {CNA};
  \node[draw, minimum width=1.7cm, minimum height=0.8cm, align=center] (sys) at (6.7,1.9) {Syst\`eme de\\refroidissement};
  \draw[->] (som.east) -- (cor.west) node[midway,above,font=\scriptsize]{$\varepsilon_k$};
  \draw[->] (cor.east) -- (cna.west) node[midway,above,font=\scriptsize]{$u_k$};
  \draw[->] (cna.east) -- (sys.west) node[midway,above,font=\scriptsize]{$u_C(t)$};
  \draw[->] (sys.east) -- (8.6,1.9) node[right,font=\scriptsize]{$T(t)$};
  % retour
  \node[draw, minimum width=1.0cm, minimum height=0.7cm] (cap) at (6.9,0.35) {Capteur};
  \node[draw, minimum width=1.6cm, minimum height=0.7cm, align=center] (far) at (4.3,0.35) {Filtre anti-\\repliement + E/B};
  \node[draw, minimum width=0.9cm, minimum height=0.7cm] (can) at (2.0,0.35) {CAN};
  \draw (8.2,1.9) node[circ]{} -- (8.2,0.35) -- (cap.east);
  \draw[->] (cap.west) -- (far.east) node[midway,above,font=\scriptsize]{$u_R(t)$};
  \draw[->] (far.west) -- (can.east);
  \draw[->] (can.west) -- (0.9,0.35) -- (som.south) node[pos=0.75,left,font=\scriptsize]{$N_k$};
\end{tikzpicture}
\end{center}

Le syst\`eme de refroidissement est mod\'elis\'e par un \textbf{premier ordre} de constante de temps $\tau = \SI{0,9}{\second}$ et de gain statique $K = \SI{4}{\degreeCelsius\per\volt}$ ; le capteur est assimil\'e \`a un gain pur $K_c = \SI{0,0316}{\volt\per\degreeCelsius}$. La r\'eponse de $u_R(t)$ \`a un \'echelon de commande $u_C(t)$ de \SI{1}{\volt} est donc celle d'un premier ordre de valeur finale $K\cdot K_c = \SI{0,1264}{\volt}$.

Le constructeur d\'esire obtenir \textbf{au moins 10 \'echantillons au cours du r\'egime transitoire}.

\begin{enumerate}[leftmargin=1.5em]
  \item \textbf{D\'eterminer} la dur\'ee du r\'egime transitoire, de $t = 0$ \`a $t_{R5\%}$. \textbf{\'Etablir} l'in\'egalit\'e que doit respecter la p\'eriode d'\'echantillonnage $T_e$, puis \textbf{en d\'eduire} la fr\'equence d'\'echantillonnage minimale $f_{min}$.
  \item \textbf{D\'efinir} le r\^ole de l'\'echantillonneur-bloqueur. \textbf{Donner} le sch\'ema \'electrique d'un \'echantillonneur-bloqueur d'ordre z\'ero.
  \item La p\'eriode d'\'echantillonnage finalement retenue est $T_2 = \SI{0,25}{\second}$. \textbf{Repr\'esenter} sur un m\^eme graphe les allures du signal $u_R(t)$, du signal \'echantillonn\'e $u_R^*(t)$ et du signal \'echantillonn\'e-bloqu\'e $u_{RB}(t)$.
  \item \textbf{D\'eterminer} le temps de conversion $T_{conv}$ maximal que doit respecter le convertisseur. \textbf{Conclure} \`a partir du document technique de l'exercice 3 sur la capacit\'e du CAN \`a convertir suffisamment rapidement.
\end{enumerate}

}{
\textbf{Corrig\'e.}

\begin{enumerate}[leftmargin=1.5em]
  \item \textbf{Dur\'ee du r\'egime transitoire.} Pour un syst\`eme du premier ordre, le temps de r\'eponse \`a \SI{5}{\percent} vaut trois constantes de temps :
  \[ t_{R5\%} = 3\,\tau = 3\times 0{,}9 = \boldsymbol{\SI{2,7}{\second}} \]
  \emph{V\'erification} : $u_R(3\tau) = 0{,}1264\times(1 - e^{-3}) = 0{,}1264\times 0{,}9502 = \SI{0,1201}{\volt}$, soit bien \SI{95}{\percent} de la valeur finale. \checkmark

  \textbf{In\'egalit\'e sur $T_e$.} Obtenir au moins 10 \'echantillons pendant le r\'egime transitoire impose que 10 p\'eriodes d'\'echantillonnage tiennent dans $t_{R5\%}$ :
  \[ 10\,T_e \leq t_{R5\%} \qquad\Longleftrightarrow\qquad \boxed{\;T_e \leq \frac{t_{R5\%}}{10} = \frac{2{,}7}{10} = \SI{0,27}{\second}\;} \]
  \[ f_{min} = \frac{1}{T_{e,max}} = \frac{1}{0{,}27} \approx \boldsymbol{\SI{3,7}{\hertz}} \]

  \emph{Attention \`a ne pas confondre les deux crit\`eres.} Celui-ci est un crit\`ere de \textbf{dynamique de boucle} : il garantit que le correcteur num\'erique « voit » suffisamment de points pour reconstruire la dynamique du proc\'ed\'e et r\'eagir avant la fin du transitoire. Le crit\`ere de Shannon, lui, porte sur le \textbf{contenu spectral} du signal \`a num\'eriser et donne une borne inf\'erieure ind\'ependante. En asservissement num\'erique, c'est presque toujours le crit\`ere de dynamique de boucle qui est \textbf{le plus contraignant} des deux -- d'o\`u la r\`egle pratique $T_e \leq T_{syst}/10$.
  \item \textbf{R\^ole de l'\'echantillonneur-bloqueur.} La conversion analogique-num\'erique n'est pas instantan\'ee : elle dure $T_{conv}$. Si la tension d'entr\'ee variait pendant cette dur\'ee, le CAN convertirait une valeur \textbf{ind\'etermin\'ee}, comprise entre celle du d\'ebut et celle de la fin de conversion. L'\'echantillonneur-bloqueur \textbf{pr\'el\`eve} la valeur du signal \`a l'instant $k\,T_e$ (phase d'\'echantillonnage) puis la \textbf{maintient constante} pendant toute la conversion (phase de blocage). Il fige ainsi l'entr\'ee du CAN et garantit que le code obtenu correspond bien \`a un instant d'\'echantillonnage pr\'ecis.

  \textbf{Sch\'ema d'un bloqueur d'ordre z\'ero} : un interrupteur command\'e (transistor MOS) charge une capacit\'e de maintien $C_H$, isol\'ee de la charge par un suiveur de tension \`a tr\`es forte imp\'edance d'entr\'ee.

  \begin{center}
  \begin{circuitikz}[european, scale=0.95, every node/.style={scale=0.95}]
    % entree + resistance serie
    \draw (0,2.6) node[left,font=\scriptsize]{$v_e(t)$} to[short, o-] (0.8,2.6);
    \draw (0.8,2.6) to[R, l=$R$] (2.4,2.6);
    % interrupteur commande K
    \draw (2.4,2.6) to[cspst, l=$K$] (4.0,2.6);
    \node[font=\scriptsize, align=center] at (3.2,3.75) {commande\\d'\'echantillonnage};
    \draw[-latex, dashed] (3.2,3.35) -- (3.2,2.95);
    % noeud de maintien
    \draw (4.0,2.6) node[circ]{} -- (4.9,2.6);
    % capacite de maintien
    \draw (4.0,2.6) to[C, l_=$C_H$] (4.0,0.7);
    \draw (4.0,0.7) -- (7.6,0.7);
    \node[ground] at (4.0,0.7) {};
    % ALI suiveur : entree + en haut, entree - bouclee sur la sortie
    \node[op amp, anchor=+, scale=0.95] (ao) at (4.9,2.6) {};
    \draw (ao.out) -- (7.9,2.6) node[right,font=\scriptsize]{$v_s(t)$};
    % contre-reaction totale (gain unite)
    \draw (ao.-) -- ++(-0.30,0) |- (6.1,1.75) -- (7.55,1.75) -- (7.55,2.6) node[circ]{};
  \end{circuitikz}
  \end{center}
  Lorsque $K$ est ferm\'e, $C_H$ se charge \`a la valeur instantan\'ee de $v_e$ (phase \emph{sample}) ; lorsque $K$ s'ouvre, $C_H$ conserve cette tension (phase \emph{hold}), le suiveur emp\^echant toute d\'echarge vers la sortie. La sortie est ainsi une suite de \textbf{paliers constants} : c'est le bloqueur d'ordre \textbf{z\'ero} (interpolation d'ordre 0).
  \item Avec $T_2 = \SI{0,25}{\second}$, les \'echantillons sont pr\'elev\'es \`a $t = 0\,;\,0{,}25\,;\,0{,}50\,;\,\dots$ s. On calcule $u_R(k\,T_2) = 0{,}1264\,(1 - e^{-k\,T_2/0{,}9})$ :

  \begin{center}
  \footnotesize
  \begin{tabular}{@{}lccccccc@{}}
  \toprule
  $t$ (s) & 0 & 0{,}25 & 0{,}50 & 0{,}75 & 1{,}00 & 1{,}50 & 2{,}00 \\
  \midrule
  $u_R$ (mV) & 0 & 30{,}7 & 53{,}9 & 71{,}5 & 84{,}8 & 102{,}5 & 112{,}7 \\
  \bottomrule
  \end{tabular}
  \end{center}

  \begin{center}
  \begin{tikzpicture}[scale=1.0, font=\footnotesize]
    % axes : 1 unite x = 1 s (echelle 2.5), y : 0.1264 V -> 2.6 cm
    \def\sx{2.45}
    \def\sy{20.0}
    \draw[->] (0,0) -- (7.9,0) node[right] {$t$ (s)};
    \draw[->] (0,-0.15) -- (0,3.25) node[above] {$u_R$ (V)};
    % graduations x
    \foreach \t in {0.5,1,1.5,2,2.5,3} {
      \draw (\t*\sx,0.06) -- (\t*\sx,-0.06) node[below, font=\scriptsize] {\t};
    }
    % graduations y
    \foreach \v/\lb in {0.0316/0{,}0316, 0.0632/0{,}0632, 0.0948/0{,}0948, 0.1264/0{,}1264} {
      \draw (0.06,\v*\sy) -- (-0.06,\v*\sy) node[left, font=\scriptsize] {\lb};
      \draw[dotted, colRule] (0,\v*\sy) -- (7.6,\v*\sy);
    }
    % signal analogique (1er ordre tau=0.9, Vf=0.1264)
    \draw[thick, colExo!45, domain=0:3.1, samples=80, smooth]
      plot (\x*\sx, {0.1264*(1-exp(-\x/0.9))*\sy});
    \node[colExo!60, font=\scriptsize, right] at (3.15*\sx,0.118*\sy) {$u_R(t)$};
    % signal bloque (escalier, Te=0.25)
    \draw[thick, colTherm]
      (0,0) -- (0.25*\sx,0)
      -- (0.25*\sx,{0.03066*\sy}) -- (0.5*\sx,{0.03066*\sy})
      -- (0.5*\sx,{0.05388*\sy}) -- (0.75*\sx,{0.05388*\sy})
      -- (0.75*\sx,{0.07147*\sy}) -- (1.0*\sx,{0.07147*\sy})
      -- (1.0*\sx,{0.08479*\sy}) -- (1.25*\sx,{0.08479*\sy})
      -- (1.25*\sx,{0.09488*\sy}) -- (1.5*\sx,{0.09488*\sy})
      -- (1.5*\sx,{0.10253*\sy}) -- (1.75*\sx,{0.10253*\sy})
      -- (1.75*\sx,{0.10832*\sy}) -- (2.0*\sx,{0.10832*\sy})
      -- (2.0*\sx,{0.11270*\sy}) -- (2.25*\sx,{0.11270*\sy})
      -- (2.25*\sx,{0.11602*\sy}) -- (2.5*\sx,{0.11602*\sy})
      -- (2.5*\sx,{0.11854*\sy}) -- (2.75*\sx,{0.11854*\sy})
      -- (2.75*\sx,{0.12045*\sy}) -- (3.0*\sx,{0.12045*\sy});
    \node[colTherm, font=\scriptsize, right] at (3.05*\sx,{0.100*\sy}) {$u_{RB}(t)$};
    % impulsions echantillonnees
    \foreach \k/\v in {1/0.03066, 2/0.05388, 3/0.07147, 4/0.08479, 5/0.09488, 6/0.10253, 7/0.10832, 8/0.11270, 9/0.11602, 10/0.11854, 11/0.12045} {
      \draw[colAuto, thick, -latex] (\k*0.25*\sx,0) -- (\k*0.25*\sx,{\v*\sy});
      \fill[colAuto] (\k*0.25*\sx,{\v*\sy}) circle (1.5pt);
    }
    \node[colAuto, font=\scriptsize] at (1.55*\sx,{0.0335*\sy}) {$u_R^*(t)$};
    % Te
    \draw[<->] (2.0*\sx,-0.62) -- (2.25*\sx,-0.62);
    \node[font=\scriptsize, right] at (2.3*\sx,-0.62) {$T_2 = \SI{0,25}{\second}$};
  \end{tikzpicture}
  \end{center}
  Le signal \'echantillonn\'e $u_R^*(t)$ (vert) est une suite d'impulsions aux instants $k\,T_2$ ; le signal \'echantillonn\'e-bloqu\'e $u_{RB}(t)$ (rouge) est l'escalier obtenu en maintenant chaque valeur jusqu'\`a l'\'echantillon suivant. On observe le \textbf{retard} introduit par le blocage : l'escalier est en permanence « en retard » d'environ $T_2/2$ sur le signal analogique, retard qui d\'egrade la marge de phase de la boucle.
  \item \textbf{Temps de conversion maximal.} Le CAN doit avoir termin\'e sa conversion avant l'arriv\'ee de l'\'echantillon suivant : la conversion doit tenir enti\`erement dans une p\'eriode d'\'echantillonnage.
  \[ T_{conv} \leq T_2 = \boldsymbol{\SI{0,25}{\second}} \]
  \textbf{Comparaison \`a la notice.} L'ADS7029-Q1 affiche un d\'ebit maximal de 2~MSPS, soit un temps de cycle (acquisition + conversion) :
  \[ t_{CYCLE} = \frac{1}{2\times 10^6} = \SI{0,5}{\micro\second} = \SI{5e-7}{\second} \]
  \[ \frac{T_2}{t_{CYCLE}} = \frac{0{,}25}{5\times 10^{-7}} = \boldsymbol{5\times 10^{5}} \]

  \smallskip
  \textbf{Conclusion.} Le CAN convertit \textbf{500\,000 fois plus vite} que ne l'exige la boucle de r\'egulation : la capacit\'e du convertisseur \`a suivre la grandeur analogique est \textbf{tr\`es largement satisfaite}, et le temps de conversion peut \^etre consid\'er\'e comme nul devant $T_2$ dans la mod\'elisation de l'asservissement. Cet \'ecart consid\'erable s'explique par la nature du proc\'ed\'e : une r\'egulation thermique est un syst\`eme \textbf{lent} ($\tau = \SI{0,9}{\second}$), alors que le composant a \'et\'e con\c{c}u pour des applications automobiles rapides. Le choix du CAN n'est donc pas ici dict\'e par sa vitesse, mais par sa r\'esolution, sa consommation et sa qualification automobile.
\end{enumerate}
}
\end{tcolorbox}

% =====================================================================
% Exercice 5 -- 2025, questions 23, 26 et 27
% =====================================================================
\begin{tcolorbox}[exoBox,
  title={Exercice 5 -- \difficulte{3} \hfill\normalfont\footnotesize 2025, \'epreuve \'ecrite disciplinaire appliqu\'ee, questions 23, 26 et 27}
]
\CAPETSHOW{
Un robot manipulateur KUKA positionne un \'ecran d'affichage sc\'enique. On \'etudie l'\textbf{axe 3} du robot :

\begin{itemize}[leftmargin=1.5em, itemsep=1pt]
  \item l'axe en sortie du servomoteur dispose d'un \textbf{d\'ebattement angulaire de $288^\circ$} ;
  \item il re\c{c}oit des consignes de position cod\'ees sur \textbf{24 bits} : le d\'ebattement complet est parcouru lorsque le mot de consigne passe de sa valeur minimale \`a sa valeur maximale ;
  \item on note $\delta\theta$ la \textbf{pr\'ecision de positionnement angulaire}, c'est-\`a-dire l'\'ecart angulaire entre deux consignes cons\'ecutives.
\end{itemize}

La position angulaire r\'eelle du moteur est mesur\'ee par un \textbf{r\'esolveur}. Un circuit de traitement convertit cette information en une \textbf{tension image} valant $\SI{-5}{\volt}$ \`a la but\'ee $\theta_{min}$ et $\SI{+5}{\volt}$ \`a la but\'ee $\theta_{max}$, de mani\`ere lin\'eaire. Cette tension est num\'eris\'ee par un CAN de plage d'entr\'ee $[\SI{-5}{\volt}\,;\,\SI{+5}{\volt}]$.

On note $\delta v$ la variation de la tension image lorsque la position angulaire du moteur tourne de l'angle $\delta\theta$.

\begin{center}
\begin{tikzpicture}[>=Stealth, scale=0.9, every node/.style={scale=0.9}, font=\scriptsize]
  \node[draw, minimum width=1.7cm, minimum height=0.85cm, align=center] (cmd) at (0,0) {Consigne\\24 bits};
  \node[draw, minimum width=1.5cm, minimum height=0.85cm, align=center] (mot) at (2.9,0) {Servomoteur\\axe 3};
  \node[draw, minimum width=1.3cm, minimum height=0.85cm, align=center] (res) at (5.5,0) {R\'esolveur};
  \node[draw, minimum width=1.7cm, minimum height=0.85cm, align=center] (tr) at (5.5,-1.7) {Traitement\\$\pm\SI{5}{\volt}$};
  \node[draw, minimum width=1.5cm, minimum height=0.85cm, align=center] (can) at (2.5,-1.7) {CAN\\$n$ bits ?};
  \draw[->] (cmd.east) -- (mot.west);
  \draw[->] (mot.east) -- (res.west) node[midway,above,font=\scriptsize]{$\theta$};
  \draw[->] (res.south) -- (tr.north);
  \draw[->] (tr.west) -- (can.east) node[midway,above,font=\scriptsize]{$\delta v$};
  \draw[->] (can.west) -- (-0.2,-1.7) -- (cmd.south) node[pos=0.8,left,font=\scriptsize]{$N$};
  \draw[->] (mot.north) -- ++(0,0.55) node[above,font=\scriptsize]{$288^\circ$ de d\'ebattement};
\end{tikzpicture}
\end{center}

Le capteur est consid\'er\'e comme \textbf{suffisamment pr\'ecis si son erreur de quantification $q$ est 10 fois plus petite} que la variation de tension $\delta v$ entre deux positions angulaires commandables du moteur.

\begin{enumerate}[leftmargin=1.5em]
  \item \textbf{Montrer} que la pr\'ecision de positionnement angulaire vaut $\delta\theta = 1{,}72\times 10^{-5}\ {}^\circ$.
  \item \textbf{\'Etablir} la valeur de $\delta v$.
  \item \textbf{D\'eterminer} le nombre de bits de codage $n$ n\'ecessaire au CAN pour que sa pr\'ecision soit suffisante. \textbf{Conclure} sur la r\'ealisabilit\'e technique de ce convertisseur.
  \item Cette conclusion sur les caract\'eristiques du servomoteur de l'axe 3 suffit-elle \`a garantir la pr\'ecision du positionnement de l'\'ecran ? \textbf{Justifier}.
\end{enumerate}

}{
\textbf{Corrig\'e.}

\begin{enumerate}[leftmargin=1.5em]
  \item \textbf{Pr\'ecision angulaire.} Le mot de consigne sur 24 bits prend $2^{24}$ valeurs ; le d\'ebattement de $288^\circ$ est d\'ecoup\'e en autant de pas :
  \[ \delta\theta = \frac{288}{2^{24}} = \frac{288}{16\,777\,216} = \boldsymbol{1{,}72\times 10^{-5}\ {}^\circ} \]
  \emph{Valeur exacte} : $1{,}716\,6\times 10^{-5}\ {}^\circ$, soit environ $\SI{0,062}{\arcsecond}$ d'arc -- une pr\'ecision de commande extr\^emement fine.

  \emph{Convention retenue} : on d\'ecoupe ici le d\'ebattement en $2^{24}$ \textbf{intervalles} (convention $2^n$), comme le fait le corrig\'e officiel du jury. Avec la convention $2^n - 1$ (nombre d'intervalles entre $2^n$ valeurs cod\'ees), on obtiendrait $288/(2^{24}-1) = 1{,}716\,6\times 10^{-5}$ : les deux r\'esultats sont \textbf{identiques aux trois chiffres significatifs} affich\'es, l'\'ecart relatif \'etant de $1/2^{24} \approx 6\times 10^{-8}$. \`A cette r\'esolution, le d\'ebat sur la convention n'a plus aucun sens physique.
  \item \textbf{Variation de tension.} La tension image est lin\'eaire sur le d\'ebattement : $\SI{10}{\volt}$ d'\'etendue ($-5$ \`a $+5$~V) pour $288^\circ$. La sensibilit\'e du capteur vaut donc $10/288 = \SI{0,03472}{\volt\per\degree}$, et pour une rotation de $\delta\theta$ :
  \[ \delta v = \frac{10}{288}\times\delta\theta = \frac{10}{288}\times\frac{288}{2^{24}} = \frac{10}{2^{24}} = \boldsymbol{\SI{5,96e-7}{\volt}} \approx \SI{0,596}{\micro\volt} \]
  \emph{Le d\'ebattement de $288^\circ$ se simplifie} : $\delta v$ ne d\'epend que de l'\'etendue de tension et du nombre de bits de la \textbf{consigne}. C'est logique : la cha\^ine de mesure doit distinguer autant de niveaux de tension qu'il existe de positions commandables.
  \item \textbf{Crit\`ere de pr\'ecision.} L'\'enonc\'e impose $q \leq \delta v/10$ :
  \[ q = \frac{\delta v}{10} = \frac{5{,}96\times 10^{-7}}{10} = \boldsymbol{\SI{5,96e-8}{\volt}} \approx \SI{6e-8}{\volt} \]
  \textbf{Nombre de bits.} La plage $[-5\,;\,+5]\,\si{\volt}$, d'\'etendue $\SI{10}{\volt}$, doit \^etre d\'ecoup\'ee en $N$ pas de largeur $q$ :
  \[ N = \frac{10}{q} = \frac{10}{5{,}96\times 10^{-8}} = 1{,}678\times 10^{8} \]
  \[ n \geq \log_2\!\left(1{,}678\times 10^{8}\right) = \frac{\ln(1{,}678\times 10^{8})}{\ln 2} = 27{,}32 \quad\Longrightarrow\quad \boldsymbol{n = 28\ \text{bits}} \]
  \emph{Arrondi \`a l'entier sup\'erieur imp\'eratif} : avec 27 bits, $2^{27} = 1{,}342\times 10^{8} < 1{,}678\times 10^{8}$, le crit\`ere ne serait pas tenu.

  \emph{Chemin de v\'erification direct} : le crit\`ere impose $2^n \geq 10\times 2^{24}$, soit $n \geq 24 + \log_2 10 = 24 + 3{,}32 = 27{,}32$. \checkmark\ On retrouve bien 28 bits, et l'on voit d'o\`u vient le r\'esultat : le facteur 10 exig\'e co\^ute $\log_2 10 \approx 3{,}32$, donc \textbf{4 bits suppl\'ementaires} par rapport aux 24 bits de la consigne.

  \smallskip
  \textbf{Conclusion sur la r\'ealisabilit\'e.} Un CAN de \textbf{28 bits} sur une plage de \SI{10}{\volt} suppose un quantum de $\SI{6e-8}{\volt} = \SI{60}{\nano\volt}$. Or ce niveau est \textbf{inf\'erieur au bruit thermique} d'une r\'esistance de quelques \si{\kilo\ohm} \`a temp\'erature ambiante sur une bande passante usuelle (quelques centaines de \si{\nano\volt} efficaces) : aucun convertisseur du commerce ne tient cette sp\'ecification en conversion rapide -- les meilleurs CAN $\Sigma\Delta$ de mesure plafonnent autour de 24 bits utiles, et encore \`a tr\`es faible d\'ebit. \textbf{L'exigence n'est donc pas r\'ealisable en l'\'etat} : c'est le crit\`ere lui-m\^eme (facteur 10 sur une consigne d\'ej\`a cod\'ee sur 24 bits) qui est irr\'ealiste, et non le choix du composant. La voie de mesure doit \^etre repens\'ee -- par exemple par un capteur incr\'emental de position \`a comptage, qui ne subit pas cette limite de dynamique analogique.
  \item \textbf{Non, cette conclusion ne suffit pas}, pour deux raisons ind\'ependantes :
  \begin{itemize}[leftmargin=1.5em, itemsep=1pt]
    \item \textbf{les autres axes} : le robot comporte 6 axes, et l'\'etude n'a port\'e que sur l'axe 3. Les erreurs de positionnement de chaque axe se \textbf{cumulent} au niveau de l'effecteur ; il faudrait v\'erifier que leur somme reste inf\'erieure \`a la tol\'erance globale ;
    \item \textbf{les autres sources d'erreur} : on n'a consid\'er\'e ici que l'erreur de \textbf{quantification} des angles de consigne. Il resterait \`a v\'erifier que la commande de chaque axe conduit \`a une \textbf{erreur statique nulle} (pr\'ecision de l'asservissement lui-m\^eme), et \`a prendre en compte les jeux m\'ecaniques, les d\'eformations \'elastiques des bras et les d\'erives thermiques.
  \end{itemize}

  \smallskip
  \textbf{Conclusion.} La v\'erification men\'ee ne valide qu'une \textbf{condition n\'ecessaire, non suffisante} : la r\'esolution de codage de l'axe 3 pourrait th\'eoriquement suffire \`a la pr\'ecision demand\'ee, mais la faisabilit\'e du CAN associ\'e n'est pas acquise et le budget d'erreur global du robot reste \`a \'etablir. C'est exactement le type de conclusion attendu au CAPET : \textbf{qualifier la port\'ee du r\'esultat} au lieu de conclure globalement sur un calcul partiel.

  \smallskip
  \textit{Point de vigilance jury (rapport de jury CAPET 2022, conseils aux candidats pour l'\'epreuve \'ecrite) : \guillemotleft\ Le jury recommande aux candidats d'apporter un soin particulier aux questions de conclusion de chacune des parties. Les \'ecarts \'evalu\'es doivent \^etre clairement mis en \'evidence et comment\'es. \guillemotright}
\end{enumerate}
}
\end{tcolorbox}

% =====================================================================
% Exercice 6 -- exercice inedit (CNA R-2R)
% =====================================================================
\begin{tcolorbox}[exoBox,
  title={Exercice 6 -- \difficulte{3} \hfill\normalfont\footnotesize exercice in\'edit}
]
\CAPETSHOW{
En sortie d'un correcteur num\'erique, la commande $u_k$ cod\'ee sur $n$ bits doit \^etre reconvertie en tension analogique. On utilise pour cela un \textbf{convertisseur num\'erique-analogique \`a r\'eseau R-2R} en configuration \emph{courant}, suivi d'un amplificateur transimp\'edance.

Le mot d'entr\'ee est $N = b_{n-1}b_{n-2}\dots b_1 b_0$ (codage binaire naturel, $b_{n-1}$ = MSB). Chaque bit $b_i$ commande un interrupteur qui aiguille le courant de sa branche soit vers la masse, soit vers l'entr\'ee inverseuse de l'ALI (masse virtuelle).

\begin{center}
\begin{circuitikz}[european, scale=0.92, every node/.style={scale=0.92}]
  % ---- trois bus horizontaux ----
  %  y = 4.3 : bus de sortie vers la masse virtuelle de l'ALI
  %  y = 3.0 : echelle R-2R
  %  y = 0.0 : bus de masse
  % ---- echelle R-2R (bus du haut) ----
  \draw (0.0,3.0) to[R, l=$R$] (1.9,3.0);
  \draw (1.9,3.0) node[circ]{} to[R, l=$R$] (3.8,3.0);
  \draw (3.8,3.0) node[circ]{} to[R, l=$R$] (5.7,3.0);
  \node[circ] at (5.7,3.0) {};
  % terminaison 2R a gauche
  \draw (0.0,3.0) node[circ]{} to[R, l_=$2R$] (0.0,0.0);
  % ---- branches 2R + interrupteurs d'aiguillage ----
  % chaque interrupteur bascule entre le bus de masse (y=0) et le bus de sortie (y=4.3)
  \foreach \x/\lb in {1.9/$b_0$, 3.8/$b_1$, 5.7/$b_{n-1}$} {
    \draw (\x,3.0) to[R, l_=$2R$] (\x,1.75);
    % contact commun
    \draw (\x,1.75) -- (\x,1.45);
    % lame de l'interrupteur (position basse = vers la masse)
    \draw[thick] (\x,1.45) -- (\x-0.42,0.75);
    \node[font=\scriptsize] at (\x+0.3,1.15) {\lb};
    % contact "masse"
    \draw (\x-0.42,0.65) node[circ]{} -- (\x-0.42,0.0);
    % contact "sortie" (renvoye vers le bus de masse virtuelle, sous le montage)
    \node[circ] at (\x+0.42,0.65) {};
    \draw (\x+0.42,0.65) -- (\x+0.42,-1.15);
  }
  % bus de renvoi vers la masse virtuelle de l'ALI (sous le bus de masse)
  \draw (2.32,-1.15) -- (9.1,-1.15) -- (9.1,5.2);
  \node[circ] at (4.22,-1.15) {};
  \node[circ] at (6.12,-1.15) {};
  % ---- bus de masse ----
  \draw (0.0,0.0) -- (7.0,0.0);
  \node[ground] at (3.35,0.0) {};
  % ---- source de reference ----
  \draw (5.7,3.0) to[R, l=$R$] (7.0,3.0);
  \draw (7.0,3.0) to[V, l=$V_{ref}$] (7.0,0.0);
  % ---- ALI transimpedance (masse virtuelle sur l'entree inverseuse) ----
  \node[op amp, anchor=-, scale=1.0] (ao) at (9.1,5.2) {};
  \draw (ao.+) -- ++(-0.5,0) node[ground, rotate=-90]{};
  \draw (ao.out) -- (11.9,5.2) node[right,font=\scriptsize]{$V_S$};
  % contre-reaction Rf
  \draw (9.1,5.2) node[circ]{} -- (9.1,6.5);
  \draw (9.1,6.5) to[R, l=$R_f$] (11.4,6.5);
  \draw (11.4,6.5) -- (11.4,5.2) node[circ]{};
\end{circuitikz}
\end{center}

On donne $V_{ref} = \SI{5,00}{\volt}$, $R = \SI{10}{\kilo\ohm}$, $R_f = R$ et $n = 12$ bits. La loi de sortie d'un tel montage est :
\[ V_S = -\,V_{ref}\,\frac{R_f}{R}\cdot\frac{N}{2^n} \]

\begin{enumerate}[leftmargin=1.5em]
  \item \textbf{Justifier} qualitativement la structure du r\'eseau R-2R : pourquoi la r\'esistance vue depuis chaque n\oe ud vers la gauche vaut-elle toujours $R$, et quelle cons\'equence cela a-t-il sur les courants de branche ?
  \item \textbf{Calculer} le quantum du CNA, ainsi que la tension de sortie pour $N = 2\,048$ puis pour $N = 4\,095$. \textbf{Commenter} l'\'ecart entre la tension obtenue \`a $N = 4\,095$ et la pleine \'echelle $V_{ref}$.
  \item Les r\'esistances du r\'eseau sont r\'ealis\'ees avec une tol\'erance relative $\epsilon$. On montre que la condition de \textbf{monotonicit\'e} du convertisseur (aucune inversion de la sortie lorsque le code cro\^it) est la plus contraignante \`a la transition du bit de poids fort, o\`u elle s'\'ecrit $\epsilon \leq 2^{-n}$. \textbf{D\'eterminer} la tol\'erance maximale admissible pour $n = 12$ bits, et \textbf{conclure} sur la technologie de fabrication \`a mettre en \oe uvre.
\end{enumerate}

}{
\textbf{Corrig\'e.}

\begin{enumerate}[leftmargin=1.5em]
  \item \textbf{Propri\'et\'e fondamentale du r\'eseau R-2R.} Tous les interrupteurs aiguillent leur branche soit vers la masse r\'eelle, soit vers la \textbf{masse virtuelle} de l'entr\'ee inverseuse de l'ALI : dans les deux cas, le bas de chaque r\'esistance $2R$ est \textbf{au potentiel z\'ero}. Le r\'eseau voit donc toujours la m\^eme topologie, quel que soit le mot d'entr\'ee.

  On raisonne alors de proche en proche depuis l'extr\'emit\'e gauche :
  \begin{itemize}[leftmargin=1.5em, itemsep=1pt]
    \item la terminaison est un $2R$ \`a la masse, en parall\`ele avec le $2R$ du premier bit : $2R\parallel 2R = R$ ;
    \item cette r\'esistance $R$ en s\'erie avec le $R$ horizontal suivant donne $2R$, de nouveau en parall\`ele avec le $2R$ du bit suivant : $2R\parallel 2R = R$ ;
    \item le motif se \textbf{reproduit \`a l'identique} \`a chaque \'etage.
  \end{itemize}
  \textbf{Cons\'equence sur les courants.} \`A chaque n\oe ud, le courant arrivant se partage entre deux branches de m\^eme r\'esistance \'equivalente ($2R$ vers le bas, $2R$ vers la gauche) : il se \textbf{divise exactement par deux}. Le courant de branche du bit $b_i$ vaut donc
  \[ I_i = \frac{V_{ref}}{2R}\cdot\frac{1}{2^{\,n-1-i}} \]
  c'est-\`a-dire une progression en \textbf{puissances de 2} strictement conforme aux poids binaires -- et ce avec \textbf{seulement deux valeurs de r\'esistance}, $R$ et $2R$. C'est tout l'int\'er\^et de la structure : un convertisseur \`a r\'esistances pond\'er\'ees ($R, 2R, 4R,\dots, 2^{n-1}R$) exigerait pour $n = 12$ un rapport de $2\,048$ entre les valeurs extr\^emes, impossible \`a appairer avec pr\'ecision sur un m\^eme substrat.
  \item \textbf{Quantum du CNA.} Avec $R_f = R$, la loi de sortie se r\'eduit \`a $V_S = -V_{ref}\,N/2^n$. Le quantum est la variation de sortie pour $\Delta N = 1$ :
  \[ q = \frac{V_{ref}}{2^n} = \frac{5{,}00}{4\,096} = \SI{1,2207e-3}{\volt} \approx \boldsymbol{\SI{1,22}{\milli\volt}} \]
  \emph{Convention} : le CNA, contrairement au CAN, ne laisse ici \textbf{aucune ambigu\"it\'e} -- le diviseur est physiquement en $2^n$, impos\'e par la structure du r\'eseau (le courant du MSB vaut exactement $V_{ref}/2R$, soit la moiti\'e du courant total possible). La convention $2^n - 1$ n'a pas lieu d'\^etre pour ce montage.

  \textbf{Tensions de sortie.}
  \[ N = 2\,048 = 2^{11} \;:\; V_S = -5{,}00\times\frac{2\,048}{4\,096} = \boldsymbol{\SI{-2,50}{\volt}} \]
  \[ N = 4\,095 \;:\; V_S = -5{,}00\times\frac{4\,095}{4\,096} = -4{,}998\,78 \approx \boldsymbol{\SI{-5,00}{\volt}} \]
  \textbf{Commentaire.} La sortie n'atteint \textbf{jamais} $-V_{ref}$ : le code maximal $4\,095$ donne $V_{ref} - q$, soit un d\'efaut de $\SI{1,22}{\milli\volt}$. C'est structurel : le mot $N = 4\,096$ n'existe pas sur 12 bits. La \textbf{dynamique r\'eelle} du CNA est donc $[0\,;\,V_{ref} - q]$, et non $[0\,;\,V_{ref}]$ -- point \`a ne pas oublier lors du dimensionnement d'un \'etage de puissance plac\'e en aval, sous peine de sous-estimer la commande maximale disponible de $1$~LSB. Le signe n\'egatif traduit simplement le montage inverseur ; un \'etage inverseur de gain $-1$ le r\'etablit si la commande doit \^etre positive.
  \item \textbf{Tolérance maximale.}
  \[ \epsilon \leq 2^{-n} = 2^{-12} = \frac{1}{4\,096} = 2{,}441\times 10^{-4} \approx \boldsymbol{\SI{0,024}{\percent}} \]
  \emph{Interpr\'etation physique de ce crit\`ere.} \`A la transition $N = 2\,047 \rightarrow 2\,048$ (passage de $0111\,1111\,1111_2$ \`a $1000\,0000\,0000_2$), \textbf{tous} les interrupteurs basculent simultan\'ement : le MSB s'active pendant que les 11 autres bits se d\'esactivent. La sortie doit cro\^itre d'exactement $1$~LSB, alors que deux contributions de taille $2\,048\,q$ et $2\,047\,q$ se compensent presque enti\`erement. Une erreur relative de $\epsilon$ sur la r\'esistance du MSB produit une erreur absolue de $\epsilon\times 2\,048\,q$ ; pour qu'elle reste inf\'erieure au pas de $1$~LSB, il faut $\epsilon\times 2^{n-1} \leq 1/2$, soit bien $\epsilon \leq 2^{-n}$. Si la condition est viol\'ee, la sortie \textbf{recule} alors que le code avance : le convertisseur devient \textbf{non monotone}, ce qui est r\'edhibitoire dans une boucle d'asservissement (inversion locale du signe du gain, donc risque d'instabilit\'e ou de cycle limite).

  \smallskip
  \textbf{Conclusion.} Une tol\'erance de \SI{0,024}{\percent} est \textbf{hors de port\'ee des r\'esistances discr\`etes} du commerce (les s\'eries courantes offrent \SI{1}{\percent}, les pr\'ecisions m\'etal \SI{0,1}{\percent} \`a co\^ut \'elev\'e, soit encore \textbf{quatre fois trop} pour ce crit\`ere). Un CNA 12 bits monotone ne peut donc \^etre r\'ealis\'e qu'en \textbf{circuit int\'egr\'e} : r\'esistances diffus\'ees sur un m\^eme substrat, \textbf{appari\'ees} par dessin (\emph{matching} obtenu par des motifs en centro\"ide commun) et ajust\'ees par \emph{laser trimming}. C'est ce qui explique qu'on n'\'etudie ce r\'eseau que comme mod\`ele de compr\'ehension et qu'on utilise en pratique un composant int\'egr\'e : ce n'est pas la loi de conversion qui est difficile \`a obtenir, mais l'\textbf{appariement} des r\'esistances. Noter enfin que seule la monotonicit\'e est ici garantie -- la pr\'ecision absolue d\'epend en outre de la stabilit\'e de $V_{ref}$ et de l'offset de l'ALI.
\end{enumerate}
}
\end{tcolorbox}

\CAPETENDFILE
